% ===== PRÉAMBULE CANONIQUE — à coller VERBATIM en tête de CHAQUE .tex =====
\documentclass[11pt,a4paper]{article}
\usepackage[utf8]{inputenc}\usepackage[T1]{fontenc}\usepackage{lmodern}
\usepackage[french]{babel}
\usepackage{amsmath,amssymb,mathtools}\usepackage{icomma}
\usepackage{siunitx}\sisetup{locale=FR}  % \qty{}{}, \si{}, \ang{}
\usepackage{xcolor}\usepackage{colortbl}
\usepackage{tikz}\usepackage{pgfplots}\pgfplotsset{compat=1.18}
\usepackage[siunitx,europeanresistors]{circuitikz} % circuits (élec.)
\usepackage[most]{tcolorbox}\usepackage{enumitem}\usepackage{array,booktabs}
\usepackage[a4paper,margin=2.2cm]{geometry}
\usetikzlibrary{arrows.meta,calc,angles,quotes,positioning,patterns,%
  backgrounds,shapes.geometric,decorations.pathmorphing,decorations.markings,intersections}
\definecolor{primary}{RGB}{222,49,99}\definecolor{primarydark}{RGB}{150,22,55}
\definecolor{defcol}{RGB}{0,114,178}\definecolor{methcol}{RGB}{52,92,148}
\definecolor{excol}{RGB}{0,135,90}\definecolor{attncol}{RGB}{200,40,55}
\tcbset{boxbase/.style={enhanced,breakable,boxrule=0.9pt,arc=2.5pt,
  left=3mm,right=3mm,top=2.3mm,bottom=2.3mm,fonttitle=\bfseries,coltitle=white}}
% --- boîtes TUTORIEL (noms EXACTS, ne pas renommer) ---
\newtcolorbox{cle}[1][]{boxbase,colback=defcol!6,colframe=defcol,
  title={\textsc{À retenir}\ifx\relax#1\relax\relax\else\textmd{ : \itshape #1}\fi}}
\newtcolorbox{methode}[1][]{boxbase,colback=methcol!7,colframe=methcol,sharp corners,
  title={\textsc{Méthode}\ifx\relax#1\relax\relax\else\textmd{ --- \itshape #1}\fi}}
\newtcolorbox{exemplebox}[1][]{boxbase,colback=excol!5,colframe=excol,
  title={\textsc{Exemple}\ifx\relax#1\relax\relax\else\textmd{ : \itshape #1}\fi}}
\newtcolorbox{piege}[1][]{boxbase,colback=attncol!6,colframe=attncol,
  title={\textsc{Piège}\ifx\relax#1\relax\relax\else\textmd{ : \itshape #1}\fi}}
\newtcolorbox{remarque}[1][]{boxbase,colback=black!4,colframe=black!45,
  title={\textsc{Remarque}\ifx\relax#1\relax\relax\else\textmd{ : \itshape #1}\fi}}
% --- boîtes EXERCICE / CORRIGÉ (noms EXACTS, auto-comptées) ---
\newtcolorbox[auto counter]{exo}[1][]{boxbase,colback=primary!4,colframe=primary,
  title={\textsc{Exercice}~\thetcbcounter\ifx\relax#1\relax\relax\else\textmd{ --- \itshape #1}\fi}}
\newtcolorbox[auto counter]{corrige}[1][]{boxbase,colback=excol!4,colframe=excol,
  title={\textsc{Corrigé}~\thetcbcounter\ifx\relax#1\relax\relax\else\textmd{ --- \itshape #1}\fi}}
\newcommand{\vv}[1]{\vec{#1}}\newcommand{\ds}{\displaystyle}
\newcommand{\R}{\mathbb{R}}\newcommand{\N}{\mathbb{N}}\newcommand{\Z}{\mathbb{Z}}
% (un \usepackage{fancyhdr}+en-tête est possible ; il est ignoré côté web)
% =========================================================================
\begin{document}
% THEME_TITLE: Composition de forces parallèles
\begin{center}
{\LARGE\bfseries\color{primarydark}Thème 2 — Composition de forces parallèles}\\[2pt]
{\color{primary}\itshape Statique — Fiche 3}
\end{center}
\vspace{1mm}

Deux forces \textbf{parallèles} $\vv F_1$ (appliquée en $A$) et $\vv F_2$ (en $B$) se remplacent par une
seule force équivalente, la \textbf{résultante} $\vv R$. Son intensité, son sens et son point
d'application $O$ dépendent du sens relatif des deux forces.

\begin{cle}[forces parallèles de même sens]
Les intensités \textbf{s'ajoutent} : $\boxed{\,R = F_1 + F_2\,}$ ; $\vv R$ est parallèle aux deux forces
et de \emph{même sens}. Le point d'application $O$ est situé \textbf{entre} $A$ et $B$, tel que
\[
\boxed{\,F_1\times\overline{AO} = F_2\times\overline{BO}\,}
\qquad\Longleftrightarrow\qquad
\frac{\overline{AO}}{\overline{BO}} = \frac{F_2}{F_1}.
\]
$O$ est toujours \textbf{plus proche de la force la plus intense}.
\end{cle}

\begin{cle}[forces parallèles de sens opposé ($F_2>F_1$)]
Les intensités se \textbf{retranchent} : $\boxed{\,F_r = F_2 - F_1\,}$ ; $\vv F_r$ a le \emph{sens de la
plus grande force}. La même règle $F_1\times\overline{AO} = F_2\times\overline{BO}$ situe $O$, mais
cette fois $O$ est \textbf{à l'extérieur} du segment $[AB]$, \emph{du côté de la plus grande force}.
\end{cle}

\begin{methode}[placer le point d'application $O$]
\begin{enumerate}[leftmargin=1.6em,topsep=1pt,itemsep=1pt]
  \item Écrire $\dfrac{\overline{AO}}{\overline{BO}} = \dfrac{F_2}{F_1}$.
  \item \textbf{Même sens} : $O$ est entre $A$ et $B$, donc $\overline{AO}+\overline{BO}=\overline{AB}$.
        \textbf{Sens opposé} : $O$ est hors de $[AB]$, donc $\overline{AO}-\overline{BO}=\pm\overline{AB}$.
  \item Résoudre le système pour obtenir $\overline{AO}$ et $\overline{BO}$.
\end{enumerate}
\end{methode}

\begin{exemplebox}[même sens, calcul complet]
$F_1=\qty{2}{N}$ en $A$, $F_2=\qty{6}{N}$ en $B$, $\overline{AB}=\qty{80}{cm}$, \emph{même sens}.
Résultante $R = 2+6 = \qty{8}{N}$. Position : $\overline{AO}/\overline{BO}=F_2/F_1=3$ et
$\overline{AO}+\overline{BO}=80$, d'où $4\,\overline{BO}=80$, soit $\overline{BO}=\qty{20}{cm}$,
$\overline{AO}=\qty{60}{cm}$. $O$ est bien plus proche de $B$ (force la plus intense).
\end{exemplebox}

\begin{piege}[sens opposé : $O$ sort du segment]
Quand les forces sont de sens opposé, $O$ n'est \emph{pas} entre $A$ et $B$ : il se place
\textbf{au-delà}, du côté de la \emph{plus grande} force (pas de la petite). Et l'on \emph{soustrait}
les intensités ($F_2-F_1$), on ne les ajoute pas.
\end{piege}

\begin{remarque}[force équilibrante]
Pour immobiliser le solide, on applique \emph{en $O$} une \textbf{force équilibrante} de même
intensité que $\vv R$ (ou $\vv F_r$) mais de \emph{sens opposé} : alors $\sum\vv F=\vv 0$.
\end{remarque}

\end{document}
